% English draft based on sections/00_introduction_ja.tex
\section{Introduction: Why Intelligence Requires Stopping}

This paper is Intelligence Part II in the VED/IFGT theory tree. Part I described intelligence as a non-closed dynamical phenomenon. Part II asks a more specific question: if inference does not close itself, how can intelligence nevertheless stop, judge, act, and persist?

The answer developed here is that intelligence is not merely the growth of inference. Intelligence is the co-evolution of inference and stopping. The inference layer $L_F$ expands justification, generates branches, and moves along gradients of closure. Yet this motion has no internal fixed point. If inference is left to inference alone, it can enter infinite regress, branching explosion, or premature pathological closure.

The central structure of this paper is therefore:
\[
L_F \;\longrightarrow\; L_R \;\longrightarrow\; L_E .
\]
The inference layer $L_F$ leaves traces. These traces form a recursive log layer $L_R$. Under external mutual constraint $H_{ij}$, the motion of $L_R$ gives rise to an emergent stopping layer $L_E$. This layer is internal to the agent in one sense, but externally dependent in another. It is not a private final answer. It is a stopping attractor made possible by histories, constraints, and shared fields.

\subsection*{Four Declarations of Position}
\addcontentsline{toc}{subsection}{Four Declarations of Position}

\textbf{First, this paper is not a theory of religion, culture, science, or society as such.}
Animism, God, religion, science, institutions, and social blockchain appear in this paper only as historical operating forms of stopping structures. The object of analysis is the emergent stopping layer $L_E$ and the ways in which finite intelligence maintains, revises, and synchronizes stopping.

\textbf{Second, religion and science are not compared as superior or inferior truth systems.}
Religion is treated as a large-scale system for the long-term maintenance of stopping structures. Science is treated as a system that institutionalizes the revision of stopping structures. The contrast is not truth versus error, but different modes of stopping operation.

\textbf{Third, stopping is not the negation of freedom.}
Without stopping, inference is not free; it is consumed by regress and explosion. $L_E$ is not introduced to suppress inference, but to make open-ended inference sustainable within finite bodies, finite time, and finite attention.

\textbf{Fourth, projection is used structurally.}
When this paper speaks of projection, it does not primarily mean psychological projection. It means a structural operation by which a state, relation, or representation acquires unequal influence within a field. In Intelligence Part II, projection appears as the projection of possible stopping loci.

\subsection{What This Paper Does Not Claim}

This paper does not claim that religion is reducible to a mistake, that science is simply a replacement for religion, or that society is a blockchain in the technical sense. It also does not prescribe how artificial intelligence should be designed or regulated. The argument is descriptive and dynamical. It identifies structural conditions under which inference can persist without closing absolutely.

\subsection{How to Read the History of Intelligence}

This paper does not read the history of intelligence as a simple history of increasing inferential capacity. The inference operator $F$ is essentially non-closed, and the increase of inferential capacity does not by itself guarantee stable intelligence. On the contrary, stronger inference can amplify branching explosion, infinite regress, and inability to stop.

The increase of inferential capacity is not desirable without qualification. Inference is operated under finite bodies, finite time, and finite attention. Inference without stopping structures becomes unsustainable. For this reason, this paper reads intelligence as the co-evolution of inference and stopping.

Animism, God, religion, science, institutions, and social blockchain are not opposed objects. They are historical operating forms for sharing, maintaining, revising, and synchronizing the emergent stopping layer $L_E$. In this sense, the history of intelligence is a history of stopping structures.

\subsection{Outline}

\S{}1 introduces the basic VED/IFGT vocabulary for the intelligence layer. \S{}2 defines the inference operator $F$. \S{}3 proves the non-closure theorem. \S{}4 analyzes responsive and silent regimes. \S{}5 introduces the three-layer dynamics $L_F/L_R/L_E$. \S{}6--\S{}9 analyze historical forms of stopping operation. \S{}10 reinterprets the Gettier problem. \S{}11 derives the instability region. \S{}12 states constructive clauses. \S{}13 concludes by marking the boundary of the theory.

\bigskip
